\documentclass[12pt]{article}

\usepackage[a4paper, total={6in, 8in}]{geometry}


\begin{document}

\textbf{CHAPTER 1
Introduction: What Is a
Dynamical System?}

\hfill

\textbf{1.1 OVERVIEW}

\hfill

Dynamical systems are all around us: from a car traveling down the road to
the ripples caused by throwing a pebble into a pond to a clock pendulum
swinging back and forth. But what is a dynamical system? First, let us
explore the words that make up the phrase “dynamical system.” The
Merriam-Webster Dictionary gives these definitions:

\hfill

\textbf{Dynamic} (adjective): always active or changing; having or showing a lot of
energy; of or relating to energy, motion, or physical force

\hfill

\textbf{System} (noun): a group of related parts that move or work together

\hfill

The term “dynamic” gives us the idea of change or motion and can deal
more specifically with physical phenomena. Mathematically, when one
thinks of change, derivatives should spring to mind; indeed, these are a key
component of how we model dynamical systems.

The definition of “system” involves groups of related parts working
together. Although this definition works in general, a diagram such as the
one shown in Figure 1.1 best illustrates the concept that we use throughout
the book. Simply stated, a \textbf{system} is an entity that has an input and an
output. A system receives an input and produces an output based on the
input and its state.

When the two words are put together to form “dynamical system,”
things get interesting. A \textbf{dynamical system} is one in which inputs, outputs,
and even the system characteristics themselves can change with time.
The relationship between input and output can be modeled mathematically
using various techniques, such as differential and difference equations, transfer functions, and state space equations. Throughout this book, we
focus on two objectives: (1) investigate various techniques and analysis
methods for dynamical systems and (2) apply these methods to examples of
real-world systems and show how they can be used in practice. As we will
see, the study of dynamical systems can reveal some very interesting
behaviors, some desirable and some not.

What about a car driving down a road makes it a dynamical system?
One example is its suspension system, represented in Figure 1.2. The input
to this system is the road, which has a varying height as the car drives down
it. The output might be the ride height of the car itself. The system consists
of the various components linking the road to the passenger’s body: the
tires, control arm, shock absorber, chassis, and so on. When you hit a bump
in the road, you experience a dynamical response: perhaps your body is
jarred by the sudden change in wheel height, or perhaps you barely notice
it. The response depends on the characteristics of the car’s suspension,
which can be modeled as a mass-spring-damper system.

\hfill

\textbf{1.1.1 Why Do We Study Dynamic Systems?}

\hfill

The main reasons are (1) to \textbf{predict} system behavior and (2) to \textbf{control}
system behavior.

An example of a system we study for prediction is the weather. Scientists,
meteorologists in particular, have spent many years and many
computing hours creating models used to predict the weather. We rely on
these predictions to plan our daily activities, from scheduling an event to
deciding what to wear. As computing power advanced and more atmospheric
data became available, the models used to predict the weather
increased in sophistication and accuracy.

The above example of the car suspension is another case of a systemic
behavior we wish to predict. Knowing the effects of variables such as spring
stiffness, fluid viscosity in shock absorbers, tire pressure, and the many other
parameters that go into the system allow the designers to predict and
therefore design a system to have the desired characteristics. Certainly,
someone buying a sports car would be disappointed if the vehicle handled
like a school bus!

There are many examples of systems whose behavior we wish to
control. The previous example of the car suspension is also one of them.
Many modern cars now have adaptive suspension systems in which the
vehicle performance through turns and over bumps in the road is monitored.
Based on this performance, the system is designed to adjust itself by
changing the damping coefficient of the shock absorbers.

Another example of controlling a system from everyday life is a thermostat
like the one shown in Figure 1.3. A thermostat allows us to set a
comfortable temperature for the room we are in, and this temperature is
maintained automatically. The temperature of the room depends on many
parameters, including heat sources within the room (e.g., heaters, lights,
people); the amount of insulation in the walls, doors, and windows; and the
temperature outside the room. Designers wishing to optimize thermostat
performance need to have an understanding of the temperature dynamics in
the room. With this knowledge, the designer can choose the controller
parameters to get the best performance from the system.

The analysis strategy for a system depends on the objective. If one is
trying to predict system behavior, then often the strategy is to model as
much of the characteristics as possible. In the case of weather, an
extraordinary amount of data is gathered for use in computer models. This
data is taken from many different sensors to measure a number of conditions
(e.g., temperature, humidity, barometric pressure) in many different
locations. Despite sophisticated modeling techniques, one aspect of weather
systems that makes prediction so difficult is its inherently chaotic behavior.
This is exemplified by the thought experiment known as the butterfly
effect: a butterfly flapping its wings in one part of the world can cause a
hurricane somewhere else weeks later. We will discuss chaos later in
Chapter 4.

If a designer is trying to analyze a system to develop a control algorithm,
the technique may not be to model as much as possible about the system. In
many cases, it is not practical and simply not necessary to obtain so much
information. For the example of a thermostat, a comprehensive model of
the room would need the temperature distribution in the space. In other
words, we would need to measure the temperature at every point in the
room. However, this kind of measurement is not practical because we
cannot possibly fit that many sensors in the room. It is more likely that we
can get away with one, two, or four sensors by making assumptions about
the temperature distribution in the room. Furthermore, the control
designer often relies on a relatively simple model and focuses on a sophisticated
control algorithm to compensate for modeling errors and
uncertainty.

\hfill

\textbf{1.2 TYPES OF SYSTEMS}

\hfill

Dynamical systems can be categorized by their characteristics. In this section,
we discuss several of these categories and give examples of each type.
Understanding the type of system you are dealing with is important because
it will drive which techniques are available for you to use. Serious problems
can occur if analysis and design methods are used on the wrong type of
system!

\hfill

\textbf{1.2.1 Continuous versus Discrete}

\hfill


Consider the above example of a car suspension system. A plot of the car’s
ride height versus time after going over a bump may look something like
the curve shown in Figure 1.4. At every instant of time, we can get a value
for the car’s height. A fraction of an instant later, we can get another value.
In fact, we can get a different value at every instant in time. This is an
example of a \textbf{continuous-time system}.

In contrast, consider money in a bank account. Assume an account is
earning 5 \% interest and there is an initial deposit of \$10,000. Figure 1.5
shows how the money would grow if no deposits or withdrawals are made
and the bank applied interest monthly. Now similar to the car suspension
example, we can check the account balance at any time and get a value.
However in this case, changes only occur once every month. Because the
in-between values are known to be held constant, we can represent the
balance as shown in Figure 1.6.

The formula for the amount of money in the account is

\[
A_n = A_0 \left (1 + \frac{r}{12} \right )^n
\]

where r is the annual interest rate, n is the number of months, A0 is the
initial deposit, and An is the amount of money in the account in month
n. This variable n represents a discrete interval of time.

The bank account is an example of a \textbf{discrete-time system}. In a
discrete-time system, the values are only recorded at particular time instants.
The in-between values do not matter or do not change (as in the bank
account), or we do not have access to them. Sampling is the process by
which a continuous-time system is converted to discrete time. In this situation,
the sampling rate is very important because if the samples are not
taken often enough, important information in the signal can be lost.1 On
the other hand, there is always an upper limit on the sampling rate caused
by hardware processing limitations, and faster sampling is not always the
key to improving system performance.

Real physical systems, such as the car suspension system, are mostly
analog in nature and therefore continuous-time systems. So why study
discrete-time systems? The major reason is that many of these systems are
computer controlled, and computers are discrete-time machines that run on
a clock. Figure 1.7 shows a typical physical system controlled by a computer.
Typically, when a computer sends signals to actuators and receives
signals from sensors, some form of analog-to-digital and digital-to-analog
conversion must take place.

\hfill

\textbf{1.2.2 Linear versus Nonlinear}

\hfill

A \textbf{linear} system can be characterized in several different, but related, ways.

\begin{itemize}
\item Its dynamics can be represented by a system of linear differential equations
(for continuous-time systems) or linear difference equations (for
discrete-time systems).
\item It has a transfer function.
\item It obeys the law of superposition.
\item A sinusoidal input produces a sinusoidal output of the same frequency
\end{itemize}

One of the most common ways to test for a system’s linearity is by
verifying if it follows the law of superposition. Superposition is usually
introduced to engineers in a circuit analysis course, but the concept applies
to linear systems in general. (Resistive circuits are linear systems after all.)
Superposition is composed of two parts, scaling and additivity.

Scaling: If a system’s input $x(t)$ results in an output $y(t)$, then an input of
$\alpha x(t)$ will result in the output $\alpha y(t)$ for any value $\alpha$. This characteristic is
shown in Figure 1.8.

In other words, if you multiply the system’s input by some value, you
end up multiplying the output by the same value. Note that this also means
zero input should result in zero output.

Additivity: If a system’s input $x_1(t)$ results in $y_1(t)$ and $x_2(t)$ results in $y_2(t)$,
then if the input is $x_1(t) + x_2(t)$, the output is $y_1(t) + y_2(t)$. A block diagram
depicting this characteristic is shown in Figure 1.9.

In other words, in a linear system, it does not matter if you sum the
signals before or after the system. The output is the same.

A \textbf{nonlinear} system is simply one that is not linear. However, there are
several reasons why a system might be nonlinear, and different classes of
nonlinearities come about because of different physical reasons.

Let’s consider an example of a simple nonlinear system. Suppose a
system has input $x$ and output $y$ and they are related by

\[
y = m x + b 
\]

where m and b are nonzero constants. This is a linear equation! So how can
it be nonlinear? Although it may be a linear equation, the system it represents
is nonlinear. We can see this if we look at the scaling property.

According to the scaling property, let $x_0(t)$ be the input and $y_0(t)$ be the
corresponding output, as in

\[
y_0 = m x_0 + b
\]

Now let’s see what happens when we multiply the input by $\alpha$ and call
the output $\gamma_1$.

\[
\gamma_1 = m (a x_0) + b
\]

And when we multiply the output by $\alpha$, we get

\[
a \gamma_0 = a (m x_0 + b)
\]

\[
=m (a x_0) + ab
\]

Because $\gamma_1$ does not equal $\alpha \gamma_0$, the system is nonlinear! We can also see
this if we put zero into the system, the output is b (not zero).

Let’s also take a look at how the system does not obey the additivity
property. (This is just an exercise. We have already shown the system is
nonlinear because it violates the scaling property.)

Suppose the inputs $x_1(t)$ and $x_2(t)$ result in outputs $\gamma_1(t)$ and $\gamma_2(t)$,
respectively. Then

\[
\gamma_1 = mx_1 + b 
\]
\[
\gamma_2 = mx_2 + b 
\]

Putting the sum of $x_1$ and $x_2$ into the system, the output is

\[
y_s = m(x_1 + x_2) + b 
\]

But

\[
\gamma_1 + \gamma_2 = mx_1 + b + mx_2 + b
\]
\[
= m(x_1 + x_2) + 2b
\]

Again we see that the system is nonlinear because $\gamma_s \neq \gamma_1 + \gamma_2$, and the
additivity property is not followed.

Here is another example that’s practical and not purely mathematical.
Consider the amplifier circuit shown in Figure 1.10.

A circuit like this is usually discussed in an introductory circuits or
electronics course. It is an inverting amplifier, and the inputeoutput
relationship is given by

\[
V_{out} = - \frac{R_2}{R_1} V_{in}
\]


If you double the input voltage, the output voltage is doubled, and the
scaling property holds. Also, if you add two inputs, the result is

\[
-\frac{R_2}{R_1} (V_1 + V_2) = \left ( - \frac{R_2}{R_1} V_1  \right ) + \left ( - \frac{R_2}{R_1} V_2  \right ) 
\]

and the additivity property holds. Clearly, this must be a linear system.

Mathematically, yes, the system is linear. However, suppose your two
input voltages are 10 and 15 V. Will the output be -25 V? No, of course
not! The circuit is only being supplied with $\pm 12 V$, and the output cannot
exceed this limit.

This op-amp circuit is an example of a system that is modeled linearly
and behaves linearly within certain limits (for input voltages less than
$\frac{R_1}{R_2} 12 V$ in magnitude). However, when you exceed those limits, the output
saturates, and the nonlinearity is seen as a distortion, as shown in
Figure 1.11. This limitation is one reason microphones clip and speakers
distort sound.
There are many types of nonlinearities, and they can come about from
different types of physical limitations such as thresholds, sticking, hysteresis,
or other imperfections. Nonlinear systems, and how to deal with them, are
discussed in more detail in Chapter 4.

\hfill

\textbf{1.2.3 Time-Invariant versus Time-Varying}

\hfill

A \textbf{time-invariant} (or \textbf{autonomous}) system is one whose behavior does
not depend explicitly on time. In other words, given an input, if you run
the system now, an hour from now, or next Tuesday, the output will be the
same. Mathematically, time invariance means that the “constants” in the
system are truly constant and do not vary with time.
As an example, consider a simple first-order model of a car:

\[
m \ddot{x} = F(t) - b \dot{x} (t)
\]

where m is the mass of the car, x is the position, F is the force applied to the
wheels by the engine and drivetrain, and b is the coefficient of friction. We
model this system as one that is time-invariant because we assume the mass
and friction to be constant.

In reality, the car is a \textbf{time-varying} (or \textbf{nonautonomous}) system. As
you drive, the car consumes fuel, and the mass decreases. However, this
change is small and occurs slowly, so mass is assumed constant. Many systems
we take to be time invariant are not because of aging. In our equations,
we use parameters such as mass, friction, length, area, resistance, and
many others as constants. But they do change as components age. The
time-invariance assumption is valid as long as the timeframe over which the
parameters change is much longer than the run time for our system. If it
takes many years for a resistor value to change, but if we run the circuit for
minutes or days, then the time-invariance assumption works. In the bank
account example, the interest rate may change, but the system can be
considered time-invariant if the interest rate is constant during the window
of time we are investigating.

To check to see if a system is time-invariant, one only needs to delay the
input and observe the output. Assuming that an input of x(t) produces an
output y(t), then the system is time invariant if the input x(t-s) gives an
output of y(t-s) for any value of s. This relationship is shown in Figure 1.12.

\hfill

\textbf{1.2.4 Memory versus Memoryless}

\hfill

A system has \textbf{memory} if its output depends at all on past inputs. A system is
\textbf{memoryless} if its output depends only on the current input. This concept
can be seen by the following example.

Consider a spring, with the relationship between position and force
given by

\[
F = k x
\]

where k is the spring constant. If you want to know the position of the
spring given its force, you simply divide by k. On the other hand, suppose
you have a system with friction where

\[
F = b \dot{x} 
\]

where b is the coefficient of friction, then given the force, you need to integrate
to get position

\[
x(t) = \frac{1}{b} \int_{- \infty}^{t} F(\tau) d \tau
\]

In this case, the force needs to be known for all time up the current time
to determine the position.

\hfill

\textbf{1.2.5 Causal versus Noncausal}

\hfill

A \textbf{causal} system is one whose output depends only on the present and the
past inputs. A \textbf{noncausal} system’s output depends on the future inputs. In a
sense, a noncausal system is just the opposite of one that has memory.

How can a real-world system be noncausal? It cannot because real
systems cannot react to the future. But noncausal systems have important
real-world applications. Consider a song stored in a sound file. Because the
entire song is stored, we could process the sound by filtering in a way that
has the current notes depend on notes later in the song. This is an example
of postprocessing in which noncausal systems may be implemented.
Another example of a noncausal system application is image processing.
The pixels to the left of the current location can be considered as the “past”
and pixels to the right as the “future.”

\hfill

\textbf{1.2.6 Deterministic versus Stochastic}

\hfill

A \textbf{deterministic} system is one in which parameters and inputs are known.
Consider the earlier example of the simple first-order model of a car:

\[
m \ddot{x} = F(t) - b \dot{x} (t)
\]

This system is deterministic if we know exactly values for the mass m,
friction coefficient b, input force F, and initial conditions of x and $\dot{x}$ and
know that there are no other inputs, or disturbances, acting on the system.
In this ideal case, we can figure out the system’s trajectory (the signals x(t)
and $\dot{x}$) simply by solving the differential equation (either analytically or
numerically).

Unfortunately, real-world systems usually have some uncertainty in
them, and figuring out the system’s trajectory is not so simple. If any one of
the parameters or inputs in the equations describing the system is not or
cannot be known exactly, then the system is stochastic. Also if a system is
affected by noise in any part of it, then the system is stochastic.

\textbf{Stochastic} systems are an active field of study involving what are known
as stochastic differential equations. In these sorts of systems, the parameters or
inputs to the system are not known exactly, but it is assumed that their
probabilistic properties are known. For example, it could be assumed that
Gaussian white noise is entering the system through the sensor data.


A more thorough discussion and exploration of stochastic systems are
beyond the scope of this book. However for interested readers, good
starting points are books by Astrom (1970) and Pugachev and Sinitsyn
(2001).

\hfill

\textbf{1.3 EXAMPLES OF DYNAMICAL SYSTEMS}

\hfill

So far, several examples were introduced in the context of the different
system types, such as the car suspension system, bank account, and op-amp
circuit. One focus of this book is to provide real-world examples for each of
the concepts discussed. The hope is to show how the theory can be applied
to real systems. As a result, many more examples will be covered
throughout the book to illustrate the concepts being discussed and then to
apply them to real-world systems. Some of these examples will carry
through multiple chapters, and many analytical methods will be applied to
them. Examples of real-world systems presented include:

\begin{itemize}
\item Driving a car, both lateral (steering) and longitudinal (speed) dynamics
\item Populations of humans and other living organisms
\item The human body, in particular the balance system
\item Aircraft engines
\item Electromechanical systems such as motor-driven devices
\item Translational and rotational mechanical systems such as the masse
springedamper and pendulum
\end{itemize}

\hfill

\textbf{1.4 A NOTE ON MATLAB AND SIMULINK}

\hfill

A second focus of this book is to illustrate analysis methods using the
common software tools of MATLAB and Simulink created by Mathworks,
Inc. As with the examples, many concepts will involve MATLAB or
Simulink (or both) to apply the methods. These programs were chosen
because of their popularity in industry and academia.

MATLAB (which stands for MATrix LABoratory) is a powerful program
that allows engineers to perform computations based on matrices
using a high-level programming language. With this program, one can
perform simulations, apply analytical tools, and display plots of results.
MATLAB has many toolboxes, which are available to expand the capabilities
into many different areas. Some examples of toolboxes relevant to
dynamical systems are the Control Systems Toolbox and System Identification
Toolbox.


Simulink is an add-on program to MATLAB that has a drag-and-drop
interface that allows users to create and simulate systems visually using block
diagrams. It has much of the same capabilities as MATLAB but with a
graphical interface, and information can be passed back and forth between
the two programs.

For this book, there is the expectation that readers have some familiarity
withMATLAB and Simulink. However, expertise is not necessary. Someone
who has used these programs should be able to follow the examples and
reproduce the results, along the way learning newelements of the software and
newstrategies for using themas analytical tools. Expertise comeswith studying
examples and practice with creating code. Readers who have not used
MATLAB or Simulink should also be able to follow the examples and
understand the results being obtained. If further help is needed, several good
books are available as MATLAB and Simulink references, including those by
Beucher andWeeks (2008), Moore (2014), and Tyagi (2012).

The MATLAB and Simulink examples show one way to address the
problem. As with any programming task, there are multiple ways to
accomplish the desired result. The provided code may or may not be the
“best” program in terms of execution time, memory use, or efficiency.
These characteristics were secondary to clarity and ease of understanding
because the main goal is conveying the ideas behind how MATLAB and
Simulink are used as a tool to solve problems.

\hfill

\textbf{REFERENCES}

\hfill

Astrom, K. (1970). Introduction to Stochastic Control Theory. Mineola, NY: Dover Publications.

Beucher, O., \& Weeks, M. (2008). Introduction to MATLAB and Simulink: A Project Approach
(3rd ed.). Hingham, MA: Infinity Science Press.

Moore, H. (2014). MATLAB for Engineers (4th ed.). Boston: Prentice Hall.

Pugachev, V., \& Sinitsyn, I. N. (2001). Stochastic Systems: Theory and Applications. River
Edge, NJ: World Scientific Publishing.

Tyagi, A. (2012).MATLAB and SIMULINKfor Engineers.NewDelhi: OxfordUniversity Press.

\hfill

\textbf{FURTHER READING}

\hfill

Coiffier, J. (2012). Fundamentals of Numerical Weather Prediction. Cambridge, UK: Cambridge
University Press.

Fayyad, S. (2012). Constructing control system for active suspension system. Contemporary
Engineering Sciences, 5(4), 189e200.

Guglielmino, E., Sireteanu, T., Stammers, C. W., Gheorghe, G., \& Giuclea, M. (2008).
Semi-active Suspension Control: Improved Vehicle Ride and Road Friendliness. London:
Springer-Verlag London.

Jackson, L. (1990). Signals, Systems, and Transforms. Reading, MA: Addison-Wesley.




\begin{thebibliography}{9}
\bibitem{texbook}
PATRICIA MELLODGE (2016) \emph{A Practical
Approach to
Dynamical Systems
for Engineers}, Woodhead Publishing Series.

\bibitem{lamport94}
Leslie Lamport (1994) \emph{\LaTeX: a document preparation system}, Addison
Wesley, Massachusetts, 2nd ed.
\end{thebibliography}

\end{document}

